\chapter{Representative Evolved Artifacts and Failure Modes}
\label{chap:artifact-appendix}
\chapterrunning{artifacts and failure modes}

This appendix presents short excerpts from archived artifacts so that the empirical claims in the
main chapters are not purely statistical. The excerpts are intentionally brief. They are included to
show the kinds of changes the loop actually produced, not to replace the aggregate evidence with
anecdotes.

\section{Successful Simple-Game Adaptation}
\label{sec:appendix-game-success}

Listing~\ref{lst:app-game-success} comes from a replay-aware pursuit/evasion run that was retained
as a functional adaptation in the phase-4 novelty audit. The main structural change is the
opponent-model helper \texttt{opp\_best\_after}, which scores each candidate move against the
opponent's likely reply rather than against the opponent's current position alone.

\begin{lstlisting}[caption={Accepted pursuit/evasion policy excerpt from the replay-aware transfer suite.},label={lst:app-game-success}]
def opp_best_after(selfx, selfy):
    opp_opts = neighbors(ox, oy)
    if not opp_opts:
        return ox, oy
    best = None
    bestv = None
    for px, py in opp_opts:
        v = d2(selfx, selfy, px, py)
        if opp_evader:
            if best is None or v > bestv:
                best, bestv = (px, py), v
        else:
            if best is None or v < bestv:
                best, bestv = (px, py), v
    return best

if self_evader:
    for nx, ny in self_moves:
        px, py = opp_best_after(nx, ny)
        v = d2(nx, ny, px, py)
        if bestv is None or v > bestv:
            bestv, best_move = v, (nx, ny)
\end{lstlisting}

\section{Novel but Unhelpful Churn}
\label{sec:appendix-game-churn}

Listing~\ref{lst:app-game-churn} is a representative novelty-spike artifact from the replay-aware
territory-control audit. It is structurally coherent and visibly different from the incumbent, but
the archived selection record marks it as rejected because its score regressed from $50.5$ to
$39.5$ against the same opponent.

\begin{lstlisting}[caption={Rejected territory-control novelty spike from the replay-aware transfer audit.},label={lst:app-game-churn}]
if (nx, ny) in ot:
    sc += 18.0
elif (nx, ny) in unq:
    sc += 10.0
elif (nx, ny) in st:
    sc += 6.0
else:
    sc += 4.0

dsc = man(nx, ny, cx, cy)
sc += -0.7 * dsc
sc += 0.35 * man(ox, oy, cx, cy)

dso = man(nx, ny, ox, oy)
sc += -0.08 * dso
\end{lstlisting}

\section{Replay-Generated TSP Heuristic Fragment}
\label{sec:appendix-tsp-heuristic}

Listing~\ref{lst:app-tsp-heuristic} shows a replay-generated TSP heuristic fragment from the
phase-5 symmetric TSPLIB suite. The evolved object is not a full solver rewrite. It is a bounded
configuration over construction, local-search, perturbation, restart, and acceptance modules.

\begin{lstlisting}[caption={Phase-5 replay-generated TSP heuristic fragment.},label={lst:app-tsp-heuristic}]
def build_heuristic():
    return {
        "construction": {
            "seed_mode": "nearest_centroid",
            "candidate_limit": 28,
            "lookahead_limit": 4,
            "distance_weight": 1.05,
            "regret_bonus": 0.45,
        },
        "local_search": {
            "two_opt_passes": 8,
            "use_three_opt": True,
            "three_opt_samples": 10,
        },
        "perturbation": {
            "enabled": True,
            "mode": "segment_reversal",
            "strength": 2,
        },
        "acceptance": {
            "mode": "annealed",
            "worse_acceptance_threshold": 0.012,
        },
    }
\end{lstlisting}

\section{Rediscovered Operator Family}
\label{sec:appendix-phase7-operator}

Listing~\ref{lst:app-phase7-operator} is a phase-7 perturbation operator that exemplifies the
double-bridge and segment-reversal family highlighted in the rediscovery analysis. The phase-7
claim is not that this operator survived as a reusable component, but that related perturbation
families reappeared repeatedly without clearing the full validation bar.

\begin{lstlisting}[caption={Phase-7 perturbation operator from the rediscovered double-bridge family.},label={lst:app-phase7-operator}]
def build_operator():
    return {
        "name": "perturbation_two_cluster_bottleneck_escape_double_bridge_compact",
        "description": "Deterministic escape schedule for 2-opt stagnation.",
        "operator_type": "perturbation",
        "primary_mode": "double_bridge",
        "secondary_mode": "segment_reversal",
        "stagnation_threshold": 3,
        "strength": 3,
        "attempts": 2,
        "escalation_strength": 2,
    }
\end{lstlisting}

\section{Bounded CVRP Solver Evolution}
\label{sec:appendix-phase9-solver}

Listing~\ref{lst:app-phase9-solver} shows a short excerpt from an accepted phase-9 CVRP solver.
The fragment is representative of the later successful solvers: a deterministic Clarke--Wright
construction remains visible, but it is followed by bounded route-improvement logic rather than by
a simple constructive pass alone.

\begin{lstlisting}[caption={Accepted phase-9 CVRP solver fragment.},label={lst:app-phase9-solver}]
def improve_swap(rs):
    if len(rs) < 2:
        return False
    best_delta = 0
    best_move = None

    for i in range(R - 1):
        ri = rs[i]
        for j in range(i + 1, R):
            rj = rs[j]
            for p in range(len(ri)):
                a = ri[p]
                for q in range(len(rj)):
                    b = rj[q]
                    if li - demands[a] + demands[b] > capacity:
                        continue
                    if lj - demands[b] + demands[a] > capacity:
                        continue
                    new_i = ri[:]
                    new_j = rj[:]
                    new_i[p] = b
                    new_j[q] = a
\end{lstlisting}

\section{Rejected Validator Case}
\label{sec:appendix-phase9-failure}

The failure modes are as important as the successes. Listing~\ref{lst:app-phase9-failure} records
a phase-9 rejection in which a large candidate was salvaged syntactically but still violated the
solver-output contract. This kind of example explains why the thesis emphasizes validator pressure
throughout the later chapters.

\begin{lstlisting}[caption={Phase-9 rejection record with validator failure.},label={lst:app-phase9-failure},language={}]
"accepted": false,
"errors": [
  "solver did not return a list of routes"
],
"generation_fallback_used": false,
"salvage_attempted": true,
"validation_issues": [
  "salvaged by trimming incomplete trailing block"
]
\end{lstlisting}
